\chapter{Bootstrap}
\label{cap:bootstrap}
% ── índice ──────────────────────────────────────────────────────
\index{bootstrap|textbf}
\index{bootstrap@\texttt{bootstrap()}|textbf}
\index{reamostragem|textbf}
\index{intervalo de confiança percentil}

% ─────────────────────────────────────────────────────────────────────────────
\section{Inferência por reamostragem}

Os erros-padrão analíticos derivados da teoria assintótica assumem, em geral,
que os resíduos seguem uma distribuição específica (normal, homocedástica,
i.i.d.). O \textbf{bootstrap} de Efron (1979) substitui premissas
distribucionais por reamostragem: replica o processo de estimação $B$ vezes
em amostras com reposição do próprio conjunto de dados.

\subsection*{Bootstrap de pares}

A abordagem mais simples reamostrea pares $(y_i, x_i)$:

\begin{enumerate}
  \item Para $b = 1,\ldots,B$: amostrar $n$ observações com reposição.
  \item Estimar $\hat\theta^{(b)}$ na amostra reamostrada.
  \item $\widehat{\mathrm{SE}}_{\mathrm{boot}} = \mathrm{sd}(\hat\theta^{(1)},\ldots,\hat\theta^{(B)})$.
\end{enumerate}

O IC percentil bootstrap de $95\%$ é o intervalo $[Q_{0{,}025},\, Q_{0{,}975}]$
da distribuição empírica de $\hat\theta^{(b)}$.

\textbf{Vantagem}: válido sem saber a distribuição de $\hat\theta$ — robusto
a erros não-gaussianos, heterocedasticidade e amostras pequenas onde a
aproximação normal é ruim.

% ─────────────────────────────────────────────────────────────────────────────
\section{Verificação por DGP}

DGP com erro uniforme em $[-2, 2]$ — não-gaussiano. Com $N = 200$, o SE
analítico do OLS pode subestimar a variância real.

\begin{lstlisting}[caption={Comparação SE analítico vs.\ bootstrap}]
seed(42)
generate df x = rnormal()
generate df e = (rnormal() - 0.5) * 4.0     // U(-2, 2) — não gaussiano
generate df y = 1.0 + 2.0 * x + e

// OLS padrão
let m_ols  = ols(y ~ x, df)
display m_ols

// Bootstrap
let m_boot = bootstrap("ols", y ~ x, df, reps=500)
display m_boot
\end{lstlisting}

\subsection*{OLS analítico}

\begin{lstlisting}[language={}, basicstyle=\ttfamily\small,
  frame=none, numbers=none, backgroundcolor=\color{codbg},
  xleftmargin=0pt, xrightmargin=0pt, breaklines=false]
OLS: R²=0.178
  const=1.046  SE=0.157   x=1.828  SE=0.279
\end{lstlisting}

\subsection*{Bootstrap ($B = 500$)}

\begin{lstlisting}[language={}, basicstyle=\ttfamily\small,
  frame=none, numbers=none, backgroundcolor=\color{codbg},
  xleftmargin=0pt, xrightmargin=0pt, breaklines=false]
════════════════════════════════════════════════════════════════════════════
                   Bootstrap SE — ols (n=500/500)
────────────────────────────────────────────────────────────────────────────
Variável    β̂       SE orig.    SE boot    IC inf      IC sup
const    1.0462     0.1566     0.1545     0.7235      1.3710
x        1.8276     0.2787     0.2679     1.3313      2.3910
════════════════════════════════════════════════════════════════════════════
\end{lstlisting}

SE bootstrap ($0{,}155$ e $0{,}268$) e analítico ($0{,}157$ e $0{,}279$)
são similares aqui, indicando que os SEs analíticos do OLS são adequados
mesmo com erros uniformes quando $N = 200$. A concordância valida a
implementação: ambos convergem para o mesmo desvio-padrão assintótico.

% ─────────────────────────────────────────────────────────────────────────────
\section{Sintaxe}

\begin{lstlisting}
// bootstrap genérico — qualquer estimador implementado
bootstrap("ols",    y ~ x, df, reps=500)
bootstrap("logit",  y ~ x, df, reps=500)
bootstrap("qreg",   y ~ x, df, reps=500, q=0.5)

// alias bootse
bootse(y ~ x, df, reps=500)                 // ols por padrão (legado)
\end{lstlisting}

\begin{center}
\begin{tabular}{lll}
\toprule
\textbf{Opção} & \textbf{Padrão} & \textbf{Descrição} \\
\midrule
\texttt{reps} / \texttt{n}  & 1000 & número de réplicas bootstrap \\
\texttt{alpha}               & 0.05 & nível para ICs (bilateral) \\
\bottomrule
\end{tabular}
\end{center}

\begin{nota}
  O estimador de bootstrap também pode estimar SEs do PSM (cap.~\ref{cap:psm}),
  onde a distribuição assintótica do ATT é não-padrão. Nesse contexto, o
  bootstrap de pares é consistente (Abadie e Imbens 2008).
\end{nota}
